\section{Learned choices of integral and propagator order}
\label{app:ordering-policy}
This appendix specifies the ordering choices and explains why consistent relabelling leaves the integrals unchanged. The comparison uses the family definitions and target patterns in Appendix~\ref{app:type-specialists}. We use parameter profile 81 for
development, profiles 82 and 83 for training, profile 84 for validation,
and profiles 85 and 86 for the reserved comparison. Each profile uses
the rational parameters defined there, while the dimension $d$ remains
symbolic. These are new parameter values of known graphs and target
patterns.

\subsection{What the choice changes}
Kira provides eight integral orderings~\cite{Kira12}. Order one treats
numerator powers as simpler than extra denominator powers, called dots;
order two reverses that preference. Orders three and four first compare
their sum and then favour numerator powers or dots, respectively.
Orders one through four compare sectors by sector number. Orders five
through eight use the corresponding index rules but first compare the
number of positive indices and favour subsectors of the declared top
sectors. We also test the four automatic propagator permutation schemes
of Kira~\cite{Kira3}. The final coefficients use the original index
convention in every comparison.

The development candidates retain the preceding method, test each integral
order from two through eight with propagator scheme one, and test
propagator schemes two through four with integral order one. These are
eleven initial choices. For each development batch, we also combine its
fastest nondefault integral order with its fastest nondefault propagator
scheme and time that pair. This adds one development calculation per
batch. The preceding seed policy remains active in each
new choice. It can use truncation, selected seed bounds, or temporary-basis
completion as defined in the preceding appendices. A changed ordering is
applied to each Kira attempt in that method. If it fails to reduce to the
prescribed basis, we repeat the preceding method with its original settings.
Both attempts enter the recorded time. This finite comparison does not
search every joint ordering or every possible elimination sequence.

To see why changing labels preserves an integral, let $\pi$ be a
permutation of the $N$ propagator positions. Set
$D'_i=D_{\pi(i)}$ and $a'_i=a_{\pi(i)}$. Reindexing the product gives
\begin{equation}
 \prod_{i=1}^{N}(D'_i)^{-a'_i}
 =\prod_{i=1}^{N}D_{\pi(i)}^{-a_{\pi(i)}}
 =\prod_{j=1}^{N}D_j^{-a_j},
 \qquad I_{D'}(\mathbf a')=I_D(\mathbf a).
 \label{eq:order-relabel}
\end{equation}
Here $I_D$ denotes the integral defined by propagators $D_i$. Since the relabelling also leaves the measure unchanged, it preserves the integral. Changing the priority of integrals instead changes the sequence in which Kira solves the IBP equations. If $A\mathbf I=0$ denotes
those equations and $P$ is the permutation matrix that reorders the integral columns, then
\begin{equation}
 \mathbf J=P\mathbf I,\qquad
 AP^{-1}\mathbf J=A\mathbf I=0.
 \label{eq:order-system}
\end{equation}
The permutation therefore represents the same equations with a different column order. In a finite calculation, however, the ordering can also affect which equations are selected, so the equivalence of representations does not replace the checks of completeness and exact output coefficients.
We require every target to reduce to the prescribed master basis and
compare every output coefficient with an independent expanded-seed
reduction. The predictor receives neither those coefficients nor the
reference reduction steps.

\subsection{Training, validation, and timing}
For each development batch, we retain the fastest candidate if its full
time is at least ten percent below the preceding method. We also retain
the preceding method. The resulting candidate set is fixed before training.
We fit the same ten predictor types per graph as in
Appendix~\ref{app:basis-completion}, using the same nine features and
complete-time log-ratio labels. Training has two measurements per
candidate and batch; validation has three. Validation chooses a fixed
control and enables a predictor only when it saves at least five percent
against that control. We freeze the models, controls, and selection requirements before
using the reserved profiles.

Each application selects its ordering before Kira creates integral weights
and generates equations. The prediction is made once per requested batch.
The selected method retains the earlier seed policy, and Kira performs all
algebraic operations. The decision selects a setting for the complete reduction before it starts. Kira continues to choose individual row operations, as required by the earlier row-model comparison.

The reserved comparison measures native choices, the preceding method,
the training-selected fixed setting for each graph, and the new selected
method. It uses five repetitions with rotating execution order. The same
worker handles all methods for each batch. The timer includes prediction,
input preparation, native equation generation, exact algebra, basis
conversion where needed, coefficient checks, fallback, and cleanup.
The caller and models are already loaded, and each attempt starts a new
Kira process. Offline development, reference construction, training, and
record writing are excluded. Additional applications receive the target
indices and master basis without reference coefficients. Their results
are checked independently after the reduction.
